\section{Evaluation is easier than generation: la palanca de la supervisión escalable}

Los métodos anteriores dependen de comparaciones humanas. Para extenderlos a modelos más capaces hay que responder una pregunta central: ¿pueden las personas evaluar tareas que ellas mismas no sabrían realizar? El punto de partida del ponente es una asimetría frecuente, aunque no absoluta: generar un resultado de alta calidad puede ser difícil, mientras que, una vez vistos los candidatos, identificar cuál es mejor suele ser más sencillo.

\subsection{De la proposición a la intuición}

La clase se detiene primero en una diapositiva de título aparte porque esta frase sostiene toda la lógica de la segunda mitad. Si la evaluation fuera tan difícil como la generation, las personas perderían todo punto de entrada a la supervisión en cuanto la capacidad del modelo las superara; si evaluar es más fácil, las personas pueden apoyarse en candidatos, evidencia y críticas para extender su capacidad de supervisión a tareas que no podrían completar por sí solas. Esta sección formula primero el principio como una ventaja de costo condicionada y después, mediante cuatro analogías, examina cuándo se cumple y cuándo deja de cumplirse porque la evidencia está oculta o porque el evaluator es manipulado.

\lecturefigure{slide-11-evaluation-title.jpg}{Principio central de la segunda mitad de la clase: evaluation is easier than generation.}{Video oficial de Stanford Online, 00:20:19--00:20:32.}

\begin{importantbox}{Formulación condicionada de la ventaja de evaluar}
Una formulación más precisa es: dados los resultados candidatos, la especificación de la tarea y evidencia suficientemente verificable, el costo de comparar o de revisar partes concretas puede ser menor que el de generar desde cero. La ventaja proviene de que los candidatos reducen el espacio de búsqueda, no de que el evaluator posea automáticamente la verdad. Si el error está oculto en un estado no visible, si la evidencia no puede obtenerse, si se puede persuadir al evaluador o si el propio objetivo de la tarea es ambiguo, la ventaja de evaluar se reduce e incluso desaparece.
\end{importantbox}

\subsection{Cómo deben leerse las cuatro analogías}

La slide enumera P/NP, ver una competencia deportiva, comprar un producto y la revisión académica por pares. Comparten una estructura: quien ejecuta necesita encontrar una solución dentro de un espacio de acciones enorme, mientras que quien evalúa, al ver unos pocos candidatos, solo necesita comparar resultados o señalar problemas concretos. Sin embargo, estos ejemplos no demuestran que «todas las tareas reales puedan verificarse con eficiencia»; solo explican por qué la preference comparison puede ser más escalable que la demonstration.

\lecturefigure{slide-12-evaluation-easier-examples.jpg}{Analogías de que evaluar es más fácil que generar: complejidad computacional, deportes, decisiones de consumo y revisión de artículos.}{Video oficial de Stanford Online, 00:20:32--00:22:25.}

\begin{knowledgebox}{Lectura de la figura: reducir las analogías a una estructura común}
P/NP subraya la diferencia entre encontrar una solución y verificar un certificado; los deportes subrayan que el público puede juzgar quién gana sin alcanzar el nivel profesional; las decisiones de consumo subrayan que comparar la experiencia de uso no exige fabricar el producto; la revisión por pares subraya que señalar defectos no equivale a escribir un artículo mejor. Lo común es que el evaluator dispone de candidatos o de evidencia. La diferencia es que la evaluación en la realidad suele tener objetivos subjetivos, variables ocultas y evidencia incompleta, por lo que no se pueden aplicar directamente las conclusiones de la teoría de la complejidad.
\end{knowledgebox}

Si el costo de generación se denota $C_{\mathrm{gen}}(T)$ y el costo de evaluación con información auxiliar $a$ se denota $C_{\mathrm{eval}}(T\mid a)$, la scalable oversight busca construir $a$ de modo que

\[
C_{\mathrm{eval}}(T\mid a) \ll C_{\mathrm{gen}}(T),
\]

donde $T$ representa la tarea y $a$ puede ser una respuesta candidata, una cita, una prueba, una critique, un diálogo o el resultado de una herramienta. El símbolo $\ll$ es un objetivo de diseño, no un teorema que se cumpla para todas las tareas.

\begin{warningbox}{El evaluador también es objeto de optimización}
En cuanto el modelo sabe qué obtiene una calificación alta, ya no optimiza solo la tarea, sino también el proceso de observación del evaluator. La expresión fluida, la evidencia seleccionada a conveniencia, la complacencia con sesgos e incluso el ocultamiento de defectos pueden elevar la calificación. Que evaluar sea más fácil solo se convierte en una señal de entrenamiento confiable cuando el propio proceso de evaluación es suficientemente robusto frente a estas estrategias.
\end{warningbox}

\subsection{Resumen de la sección}

La ventaja de evaluar ofrece un punto de entrada a la scalable oversight: las personas no necesitan generar por sí solas resultados de nivel experto y aun así pueden comparar candidatos o revisar evidencia concreta. Pero esta ventaja depende de evidencia visible, especificaciones claras y un proceso de evaluación resistente a la manipulación; no es una garantía de que «las personas siempre podrán supervisar a una IA más capaz».
